\section{Hear the Warning from Christ}

In the Sermon on the Mount, the warning about hell enters the ordinary relations of human life. Christ speaks of anger against a brother and contemptuous speech, then commands reconciliation before the offering is brought to the altar (Mt 5:21--24). Worship and conduct toward another person cannot be sealed off from one another. A mouth that calls upon God must also answer for the words by which it humiliates someone God has made. An injury carefully preserved in memory can become the excuse for further injuries, until cruelty seems like justice to the person practicing it.

Christ addresses the heart from which such conduct proceeds. The commandment against killing reaches our treatment of the brother; the commandment against adultery reaches the deliberate lust of the heart (Mt 5:27--30). The outward act matters, and the interior consent matters. The Lord sees what can remain hidden beneath a respectable life. There are sins for which no public accusation will ever come, but concealment does not make them harmless or remove them from his judgment.

His words about the eye and the hand make the cost of fidelity unmistakable. Something as dear as a bodily member cannot be preferred to eternal life. The command calls for the decisive rejection of what leads into sin. St.~Augustine, reading the parallel warning in Mark, applies the members to people held as close and useful as parts of one's body: even such an attachment must be relinquished when it draws one away from God (\emph{City of God} XXI.9). The renunciation can be painful precisely because the attachment is strong.

Christ's command concerns the rejection of sin and its occasions; bodily mutilation would injure the person called to salvation. Augustine's example makes the severity concrete. An attachment may give companionship, practical help, and the comfort of being understood while also drawing a person repeatedly into grave sin. Those real benefits explain why relinquishing it hurts. They do not remove the obligation to end the sinful conduct. Honest changes may preserve the relationship's good; where the relationship itself remains an occasion of grave sin, ending it may be required. Obedience then involves accepting a real loss and seeking God's help to endure it.

The Sermon ends by bringing religious profession under the same light. The narrow gate leads to life; the broad road leads to destruction (Mt 7:13--14). The breadth of a road says nothing about the goodness of its destination. Common practice does not absolve a conscience. If a dishonesty has become ordinary within a profession, or a contempt has become normal within a household, the number of people sharing it cannot make it faithful to Christ.

The following verses speak of fruit, of calling Jesus Lord, and of doing the Father's will. Christ even describes people appealing to extraordinary works performed in his name, yet hearing that they must depart because they have worked iniquity (Mt 7:15--23). Religious language and religious gifts do not excuse an unrepented evil life. We may know the prayers, defend the faith, or be admired for a visible service and still need to obey the command we have evaded.

The two houses complete the passage. Both builders hear; one puts the words into practice, while the other does not (Mt 7:24--27). The difference appears when the houses are tried. A meditation likewise can leave the hearer convinced while a known duty remains undone. Someone who has withheld a debt may understand repentance more clearly after prayer and still avoid the creditor. Paying what is owed requires an act beyond understanding its necessity. The Lord compares the security of the house to that obedience, so the hearer must consider where his conduct still refuses what his mind acknowledges.

In Matthew's parable of the wheat and the cockle, the master permits both to grow until the harvest (Mt 13:24--30). The present mingling of good and evil will not last forever. The servants do not determine the hour, and their eagerness to uproot everything is restrained by concern for the wheat. When Christ explains the parable, he identifies the harvest with the end of the world and the reapers with the angels. The Son of man sends them; they gather the workers of iniquity and cast them into the furnace of fire, while the just shine in their Father's kingdom (13:36--43).

God restrains the servants while the wheat and cockle grow together, and Christ assigns the final separation to his angels. His patience gives sinners time to repent. It also forbids us to treat our own anger as the execution of his sentence. St.~Alphonsus begins his consideration on the abuse of mercy with this parable: the Lord who waits for the sinner also judges persistent refusal (\emph{Preparation for Death}, XVII.1). The sinner misuses that patience when he takes continued prosperity as permission to continue in sin.

Matthew places the parables of the treasure and the pearl between this explanation and the parable of the net (13:44--50). The kingdom is worth surrendering everything to possess. The net then presents the final separation and repeats the warning about the furnace. Desire for the kingdom and fear of its loss thus stand close together in Christ's teaching. We cannot understand what it means to lose heaven while treating heaven as one modest good among many. The good offered is so great that every other possession must be judged in relation to it.

In Mark, the warning follows Christ's care for the little ones and his promise concerning a cup of water given to those who belong to him (Mk 9:40--47, in the Douay--Rheims numbering). The good done to another is precious; leading another into sin is grave. The hand, foot, and eye recur, together with the worm that does not die and the unquenchable fire. Christ impresses the truth by repetition. Love must refuse to make another person's fall the price of our pleasure, approval, or profit. We answer for the influence we willingly exert as well as for the acts we perform alone.

Lazarus lies at the gate while the rich man feasts daily. The need is nearby, embodied, and persistent. Death comes to both, and their conditions are reversed: Lazarus is comforted, while the rich man suffers torment and begs for relief (Lk 16:19--26). The story follows the teaching that one cannot serve God and mammon, and Christ's rebuke to those who justify themselves before others while God knows their hearts (16:13--15).

The great chasm in the parable cannot be crossed. The rich man then asks that his brothers be warned, but Abraham directs them to Moses and the prophets. Their duty is to hear the word already given; an extraordinary sign would not supply the obedience they refuse to that word (16:27--31). The living reader has received the warning the brothers need. There is no need to wait for an unusual experience before attending to a neighbor or repenting of a known sin.

Matthew's final judgment scene makes the encounter with Christ still more direct. The King names food, drink, welcome, clothing, care, and a visit. Those condemned answer as though they had never met the Lord in need, but he identifies their omission toward the least with their omission toward himself (Mt 25:31--46). The Catechism receives this warning as a summons to meet the serious needs of the poor and the little ones (CCC 1033). Charity has bodily tasks. A person can refuse it through settled indifference as well as through active hostility.

Christ knows the life that public reputation can conceal. His words about the builders and those who call him Lord give religious hearers particular cause to examine themselves. Familiarity with his teaching can make its language easy to repeat while obedience remains costly. The warning reaches us as people who have heard, and who may therefore mistake hearing for fidelity. In annual recollection we must allow a command we already know to question the conduct we have continued to defend. The Judge who will expose that refusal is speaking while it can still be confessed and forsaken.

